\chapter{Software para modelizaci\'{o}n, análisis, procesamiento y síntesis de se\~{n}ales}

\section{Introducci\'{o}n a Matlab / Octave}

Matlab y Octave son int\'{e}rpretes que nos permite realizar c\'{a}lculos
num\'{e}ricos, incluso con n\'{u}meros complejos, vectores y matrices,
y realizar gr\'{a}ficos en distintas representaciones.

En este apunte, mostramos algunas de las operaciones y comandos b\'{a}sicos para
la operaci\'{o}n interactiva. 

\subsection{La l\'{\i}nea de comandos}

\begin{enumerate}

\item Cuando Matlab/Octave comienza a ejecutarse, aparece una ventana tipo consola 
 (sin interfaz gr\'{a}fica), y una l\'{\i}nea marcada con el s\'{\i}mbolo
  \verb->>-. Esta es la l\'{\i}nea de comandos, donde escribimos las
  diferentes ordenes.

\item Para probar podemos escribir
\begin{verbatim}
2 + 2
3 - 5
3 / 1.4
2 * pi / 4
(1 + 2 * i) / (1 - 2 * i)
\end{verbatim}
Nota: $i$ es la unidad imaginaria.

\item En las \'{u}ltimas versiones de Matlab/Octave, a medida que vamos
  escribiendo comandos, estos quedan registrados a una historia de
  comandos, que podemos revisar m\'{a}s tarde y recordar lo hecho.

\item El comando diary nos permite grabar en un archivo de texto todas las
ordenes, datos y resultados de una sesi\'{o}n. Para usar el comando
diary escribamos en la l\'{\i}nea de comando
\begin{verbatim}
>> diary 'archivo.txt'
\end{verbatim}

\item La idea de esta clase es generar un archivo, denominado
  \verb-archivo.txt- con todas las \'{o}rdenes dadas.

\end{enumerate}

\subsection{El sistema de ayuda}

\begin{enumerate}
\item Para ver como se usa un comando es posible escribir en la l\'{\i}nea
de comando
\begin{verbatim}
>> help comando
\end{verbatim}
donde comando es el nombre del comando que nos interesa.

\item La orden help comando nos da una descripci\'{o}n del comando y
posiblemente una lista de comandos relacionados.
\begin{verbatim}
>> help diary
\end{verbatim}
\end{enumerate}

\subsection{Valores escalares}

\begin{enumerate}
\item Como hemos visto, operamos con n\'{u}meros reales y complejos,
tambi\'{e}n llamados escalares.

\item La unidad imaginaria puede escribirse $i$ o $j$.

\item Probemos los siguientes comandos
\begin{verbatim}
(1 + 2 * j) * (1 - 2 * i)
45 * (56 + i) - 34 * (16 - 32 * j)
\end{verbatim}
\end{enumerate}

\subsection{Valores matriciales o arreglos}

\begin{enumerate}
\item Matlab/Octave trabaja tambi\'{e}n con arreglos de reales y
  complejos. 

\item Podemos crear un vector fila escribiendo
\begin{verbatim}
[ 1 + 2 * i 1 - 3 * i]
\end{verbatim}
\begin{verbatim}
[ 1 + 2 * i, 1 - 3 * i]
\end{verbatim}

\item Podemos crear un vector columna escribiendo
\begin{verbatim}
[ 1 + 2 * i; 1 - 3 * i]
\end{verbatim}
o
\begin{verbatim}
[ 1 + 2 * i
  1 - 3 * i]
\end{verbatim}
\end{enumerate}

\subsection{Met\'{a}fora de Matlab/Octave}

Podemos suponer, a modo de met\'{a}fora, que disponemos de un
pizarr\'{o}n donde escribimos las cuentas que vamos haciendo,
l\'{\i}nea por l\'{\i}nea.

\subsection{Variables}

\begin{enumerate}

\item A diferencia de una calculadora, Matlab/Octave permite asignar nombres
  a los diferentes valores. Esta combinaci\'{o}n de nombre y valor
  recibe el nombre de variable.

\item A continuaci\'{o}n vemos ejemplos de los distintos tipos de variables.
La constante $i$ es la unidad imaginaria, y ya est\'{a} definida en
Matlab. El s\'{\i}mbolo \% se usa para indicar que aquello escrito a
su derecha hasta el fin de la l\'{\i}nea es un comentario

\item Pruebe ingresar lo que aparece entre \verb->>- y \verb-%-
\begin{verbatim}
>> a = 5        % escalar real
>> b = 5 + 10 * i   % escalar complejo.
>> c = [1 2 3]   % vector fila
>> d = [1, 2, 3]    % otra forma de escribir un vector fila
>> e = [1; 2; 3]    % un vector columna
>> f = [1, 2, 3; 4, 5, 6; 7, 8, 9]; % una matriz 3 x 3
\end{verbatim}

\item Matlab considera comentario todo lo que aparece desp\'{u}es del s\'{\i}mbolo
  \verb-%-, y no lo ejecuta.

\end{enumerate}

\subsection{Continuando con la met\'{a}fora de Matlab / Octave}

\begin{enumerate}
\item Las variables son similares a divisiones que dibujamos en nuestro
pizarr\'{o}n, en las cuales escribimos los valores y los nombres
respectivos.

\item El espacio ocupado por las variables se denomina espacio de
  trabajo.
\end{enumerate}

\subsubsection{Los comandos who y what}

\begin{enumerate}
\item El comando who da una lista de las variables presentes en el espacio
de trabajo (workspace)

\item Pruebe ingresar
\begin{verbatim}
>> help who
\end{verbatim}

\item El comando what da una lista de archivos matlab presentes en un
directorio.  Verifique cual es el directorio por defecto.

\item Pruebe ingresar
\begin{verbatim}
>> help what
\end{verbatim}

\item Pruebe ingresar
\begin{verbatim}
>> who
>> what
\end{verbatim}

\end{enumerate}

\subsection{Operaciones con n\'{u}meros complejos}

\begin{enumerate}
\item A continuaci\'{o}n vemos algunas de las operaciones que pueden
realizarse con n\'{u}meros complejos

\item Pruebe ingresar lo que aparece entre \verb->>- y \verb-%-
\begin{verbatim}
>> x = 3 + 4 * i
>> real(x)  % 3
>> imag(x)  % 4
>> conj(x)  % 3 - 4 * i
>> abs(x)   % 5
>> angle(x) % 0.9273
\end{verbatim}
\end{enumerate}

\subsection{Ra\'{\i}ces y coeficientes de un polinomio}

\begin{enumerate}
\item Supongamos que tenemos un polinomio
\begin{equation}
  b_{n}x^{n}+b_{n-1}x^{n-1}+...+b_{0}=0.
\end{equation}
Para hallar las ra\'{\i}ces del polin\'{o}mio podemos colocar los
coeficientes del polinomio en un vector y usar la funci\'{o}n $roots$
\begin{verbatim}
p = [bn, ..., b0];
roots(p)
\end{verbatim}

\item Supongamos que tenemos que hallar las ra\'{\i}ces de un polinomio
\begin{equation}
  x^{2}+2x+3=0.
\end{equation}
Escribamos entonces
\begin{verbatim}
p = [1,2,3];
roots(p)
\end{verbatim}

\item Supongamos ahora que tenemos las ra\'{\i}ces y necesitamos hallar los
coeficientes del polinomio. Podemos escribir entonces
\begin{verbatim}
>> r = [-1.0, 2.7, -3.1, -3.1];
>> poly(r)
\end{verbatim}

\item Por ejemplo, si las ra\'{\i}ces son
\begin{equation}
  1,-1+2i,-1-2i,
\end{equation}
entonces para hallar los coeficientes del polinomio, debemos escribir
\begin{verbatim}
>> r = [1,-1+2*i,-1-2*i];
>> poly(r)
\end{verbatim}

\item Pruebe con
\begin{verbatim}
>> p = [1, 2, 3, 4, 5];
% el vector correspondiente al polinomio x^4+2x^3+3x^2+4x+5
>> r = roots(p)
>> poly(r)
\end{verbatim}

\item \textquestiondown Cu\'{a}l le parece que es el resultado de la
siguiente orden?
\begin{verbatim}
poly(roots(p))
\end{verbatim}

\item Pruebe las siguientes \'{o}rdenes
\begin{verbatim}
>> help roots
>> help poly
\end{verbatim}

\end{enumerate}

\subsection{Generaci\'{o}n de vectores y matrices}

\begin{enumerate}
\item Matlab permite la generaci\'{o}n de valores y su asignaci\'{o}n a
vectores y matrices mediante un mecanismo sumamente eficiente. 

\item Por ejemplo, supongamos que deseamos obtener un vector con una lista de
n\'{u}meros que comienza en $1,$ termina en $2$ con
incrementos de $0.1$. Escribiendo
\begin{verbatim}
>> t = 1:0.1:2
\end{verbatim}
se obtiene el vector deseado. Observe que $t$ es un vector fila, es
decir, est\'{a} formado por una fila de once columnas ($1 \times 11).$

\item Si deseamos generar $t$ como un vector columna debemos escribir
\begin{verbatim}
>> t = (1:0.1:2)'
\end{verbatim}

\item Practique los siguientes puntos
\begin{enumerate}
\item Genere un vector que contenga valores entre $-2$ y $2$ separados
  $0.01$.

\item En especial, para generar matrices llenas de ceros o unos
  podemos usar las funciones $Zeros$ y $Ones$, respectivamente.
  Encuentre como se usan estas funciones y haga un ejemplo con cada
  una.
\end{enumerate}

\end{enumerate}

\subsection{Operaciones sobre vectores y matrices,  primera parte}

\begin{enumerate}
\item Matlab permite realizar operaciones con matrices siempre que sus
dimensiones sean compatibles. Por lo tanto, primero veremos cu\'{a}l es
la funci\'{o}n que nos permite averiguar las medidas de una variable.

\item Pruebe los siguientes comandos
\begin{verbatim}
>> t = 1:0.1:2;
>> size(t)
>> size(t')
\end{verbatim}

\item Las operaciones de suma, resta y multiplicaci\'{o}n por escalares son
muy simples. Pruebe los siguientes comandos
\begin{verbatim}
>> t + 2
>> t - 2
>> t * 2
\end{verbatim}

\item Estas operaciones se realizan sumando $2$, restando $2$ y
multiplicando por $2$ cada elemento del vector. De la misma forma se
pueden aplicar funciones trigonom\'{e}tricas
\begin{verbatim}
>> cos(t)
>> sin(2 * t - 1)
\end{verbatim}

\item Supongamos ahora que queremos elevar al cuadrado cada elemento de un
vector. Pruebe con los comandos
\begin{verbatim}
>> t .* t
>> t .^2
\end{verbatim}

\item Para realizar el producto de dos vectores escriba
\begin{verbatim}
>> t' * t
>> t * t'
\end{verbatim}
Indique la diferencia entre ambas operaciones

\end{enumerate}

\subsection{Operaciones sobre vectores y matrices, segunda parte}

Es posible realizar operaciones vectoriales como el producto o la
exponenciaci\'{o}n elemento a elemento, como la suma o la resta.

\begin{enumerate}
\item Definamos
\begin{verbatim}
x = [-1 2]
y = [3 4]
\end{verbatim}

\item Si calculamos $z = x + y$ resulta $z = [(-1 + 3) , (2+ 4)]$.

\item Si calculamos $z = x * y$ obtenemos un mensaje de error,
  diciendo que los vectores no tienen dimensiones adecuadas para ser
  multiplicados. \textquestiondown C\'{o}mo podemos obtener 
  $[(-1 * 3), (2 * 4)]$?.

\item Es posible emplear el operador $.*$ para multiplicar elemento a
  elemento. Pruebe
\begin{verbatim}
x .* y
\end{verbatim}
y obtendr\'{a}
\begin{verbatim}
[-3 8]
\end{verbatim}

\item Lo mismo se puede hacer con la divisi\'{o}n 
\begin{verbatim}
x ./ y
\end{verbatim}
y la exponenciaci\'{o}n
\begin{verbatim}
x .^ y
\end{verbatim}
o
\begin{verbatim}
x .^ 2
\end{verbatim}

\end{enumerate}

\subsection{Operaciones sobre vectores y matrices, tercera parte}

\begin{enumerate}

\item Ingresemos un vector
\begin{verbatim}
v = [1, 2, 3, 4, 5, 6, 7]
\end{verbatim}

\item Podemos manipular los valores almacenados en un vector $v$ empleando
  la notaci\'{o}n \verb-v(i)-, donde $i$ es un n\'{u}mero entero mayor
  que $0$. Pruebe
\begin{verbatim}
v(1)
\end{verbatim}

\item Tambi\'{e}n podemos manipular una parte del vector $v$
  escribiendo \verb-v(i:j)-, donde $i$ es un n\'{u}mero entero mayor
  que $0$, y $i <= j$. Pruebe
\begin{verbatim}
v(2:7)
v(1:6)
v(2:7) - v(1:6)
\end{verbatim}

\item Observe que tambi\'{e}n es posible escribir
\begin{verbatim}
v(2:end) - v(1:end-1)
\end{verbatim}

\item Matlab/Octave tiene una extensa librer\'{\i}a de funciones aplicables a
vectores y matrices. A continuaci\'{o}n mostramos una peque\~{n}a
tabla representativa
\begin{center}
  \begin{tabular}{ll}
    Comando & Resultado \\
    det & determinante de una matriz \\
    eig & autovalores y autovectores \\
    inv & inversa matricial \\
    svd & valores singulares de una matriz \\
    cond & numero de condicion de una matriz
  \end{tabular}
\end{center}

\item Supongamos un sistema de $2$ ecuaciones con dos inc\'{o}gnitas
\begin{equation}
  Ax=b
\end{equation}
donde 
\begin{verbatim}
  A = [ 1, 2; 2, 1];
  b = [ 1; 2];
\end{verbatim}
\begin{enumerate}
\item Como se puede resolver este sistema de ecuaciones con alguna de
  las funciones indicadas arriba?
\begin{verbatim}
  x = inv(A) * b;
\end{verbatim}

\item Como puede asegurar que el resultado es correcto?
\begin{verbatim}
  x = inv(A) * b;
  e = A * x - b;
  e' * e   % e' es el vector transpuesto
\end{verbatim}

\item Resuelva los sistemas de ecuaciones con los siguientes
  coeficientes y calcule los errores correspondientes
  \begin{eqnarray}
    A_{1} &=&\left[
      \begin{array}{cc}
        1 & 1 \\
        -1 & 1
      \end{array}
    \right] ,b_{1}=\left[
      \begin{array}{c}
        0 \\
        1
      \end{array}
    \right] , \\
    A_{2} &=&\left[
      \begin{array}{cc}
        100.0 & 100.0                               \\
        100.0 & 100.0 + 5 * eps
      \end{array}
    \right] ,b_{2}=\left[
      \begin{array}{c}
        0 \\
        1
      \end{array}
    \right] .
  \end{eqnarray}

\item \textquestiondown Para que sirve la variable $eps$?

\item Escriba un peque\~{n}o resumen del uso de las operaciones
  matriciales empleadas hasta ahora junto con sus conclusiones.

\item Busque en la ayuda de Matlab/Octave el listado de operaciones
  matriciales disponibles.
\end{enumerate}

\end{enumerate}

\subsection{Graficaci\'{o}n de funciones}

\begin{enumerate}
\item Para graficar funciones existen en Matlab/Octave distintos comandos. 

\item Escriba
\begin{verbatim}
>> t = 0:0.1: 2* pi;
>> y = cos(2 * t - 1);
>> figure(1); plot(t,y);
>> figure(2); stem(t,y);
\end{verbatim}

\item Como ve, $figure$ sirve para generar distintas ventanas de
graficaci\'{o}n, $%
plot$ reconstruye una funci\'{o}n continua a partir de una secuencia de
valores discretos uniendo los puntos con rectas, y $stem$ grafica una
secuencia de valores discretos.

\item Para obtener su propia interpretaci\'{o}n de estos comandos escriba
\begin{verbatim}
>> help figure
>> help plot
>> help stem
\end{verbatim}

\item Como puede ver, estos comandos tienen muchas opciones y est\'{a}n
relacionadas con muchos otros comandos. Por ejemplo, si desea generar
distintos distintos tipos de l\'{\i}neas y colores para los
gr\'{a}ficos puede escribir \verb-plot (t, y, s)- donde $s$ es un
string o cadena de car\'{a}cteres del tipo \verb-abc- donde $a$ es
cualquiera de los siguientes c\'{o}digos
\begin{verbatim}
y (yellow) amarillo, 
m (magenta o violeta), 
c (cyan o celeste), 
r (red o rojo), 
g (green) verde, b (blue)
azul, w (white), k (black) negro
\end{verbatim}
$b$ es cualquiera de los siguientes c\'{o}digos
\begin{verbatim}
. punto, 
o circulo, 
x equis, 
* asterisco, 
s (square) cuadrado, 
d diamante o rombo, 
v la letra v (abajo), 
< menor (izquierda), 
> mayor (derecha),
^ carat (arriba), 
p pentagrama, 
h hexagrama
\end{verbatim}
y $c$ es cualquiera de los siguientes c\'{o}digos
\begin{verbatim}
- solida, 
: punteada, 
-. barra y punto, 
-- barra.
\end{verbatim}

\item Por ejemplo,
\begin{verbatim}
>> plot(t,y,'c+:')
\end{verbatim}
dibuja una l\'{\i}nea color celeste con la interpolaci\'{o}n punteada
uniendo las cruces en cada dato.

\item Observe que con $grid$ puede superponer una grilla al dibujo, con $%
xlabel$ coloca un t\'{\i}tulo al eje de coordenadas, $ylabel$ coloca un
t\'{\i}tulo al eje de absisas, $title$ coloca un t\'{\i}tulo al gr\'{a}fico y $%
clf$ borra el gr\'{a}fico anterior si lo hubiera. Pruebe

\begin{verbatim}
>> t = 0:0.1:2*pi;
>> clf; plot(t,cos(t),'m:o');
>> grid;
>> xlabel('t');
>> ylabel('cos(t)');
>> title('funci\'{o}n peri\'{o}dica'); % Latex para acentos, solo en Matlab
\end{verbatim}

\item Un comando que nos resultar\'{a} \'{u}til es $subplot.$ Nos permite
componer distintos subgr\'{a}ficos en una misma ventana. Pruebe
\begin{verbatim}
>> clf;
>> subplot(3,2,1); plot(t,cos(t));
>> subplot(3,2,2); stem(t,cos(t));
>> subplot(3,2,3); plot(t,sin(t));
>> subplot(3,2,4); stem(t,sin(t));
>> subplot(3,2,5); plot(t,tan(t));
>> subplot(3,2,6); stem(t,tan(t));
>> help subplot
\end{verbatim}

\end{enumerate}

\subsection{Como guardar gr\'{a}ficos en archivos postcript, Matlab}

\begin{enumerate}

\item Supongamos que hemos realizado un gr\'{a}fico cualquiera y lo queremos
guardar en un archivo para incluirlo despu\'{e}s en un informe.

\item Pruebe
\begin{verbatim}
>> t = 0:0.1:2*pi;
>> clf; plot(t,cos(t));
>> print -deps cos.eps
\end{verbatim}

\item Vea cuales son las opciones del comando print con el sistema de
  ayuda.

\item El archivo generado $cos.eps$ puede incluirse en un documento
  Latex, y a continuaci\'{o}n generar un documento pdf.

\end{enumerate}

\subsection{Como guardar gr\'{a}ficos en archivos postcript, Octave}

\begin{enumerate}
\item Pruebe
\begin{verbatim}
>> t = 0:0.1:2*pi;
>> clf; plot(t,cos(t));
>> print('my_plot.tex', '-dtex');
\end{verbatim}
\end{enumerate}

\subsection{Graficaci\'{o}n de funciones en tres dimensiones}

\begin{enumerate}
\item Es posible graficar curvas y superficies en ``tres dimensiones''. Por
ejemplo, podemos usar el comando plot3.

\item Busque la ayuda de plot3.
        
\item Haga un gr\'{a}fico con plot3.

\item Grabe el gr\'{a}fico en un archivo.
\end{enumerate}

\subsection{Escritura de texto, o salida de datos en forma de texto}

\begin{enumerate}
\item Con el comando fprintf podemos escribir texto en la pantalla
  (salida estandar).

\item Busque la ayuda de fprintf.

\item Compare el resultado de
\begin{verbatim}
>> fprintf('hola')
\end{verbatim}
con
\begin{verbatim}
>> fprintf('hola\n')
\end{verbatim}
El caracter \verb-\n- indica que el cursor debe moverse a la
l\'{\i}nea siguiente.

\item Es posible escribir datos de una forma similar a la de un
  programa en C
\begin{verbatim}
>>  fprintf('entero %d real %f\n', 1, 1.44);
>>  fprintf('strings %s\n', 'cadenas');
\end{verbatim}

\end{enumerate}

\subsection{Iteraci\'{o}n de comandos}

\begin{enumerate}

\item Supongamos que deseamos armar una tabla de la funci\'{o}n $y = 2
  \, x - 1$ para todos los valores enteros de $x$ entre $1$ y $8$.

\item Podemos usar el comando 
\begin{verbatim}
for x = 1:8
  fprintf('x = %d; y = %d\n', x, 2 * x - 1);
end
\end{verbatim}

\item Podemos usar el comando 
\begin{verbatim}
x = 1;
while (x <= 8)
  fprintf('x = %d; y = %d\n', x, 2 * x - 1);
  x = x + 1; 
end
\end{verbatim}

\end{enumerate}

\subsection{Toma de decisiones}

\begin{enumerate}

\item la funci\'{o}n \verb-mod(x, y)- calcula el resto de la
  divisi\'{o}n entera, o sea el n\'{u}mero $r$ tal
  que $x = q * y + r$, donde $y > r >= 0$, donde $q$, $r$, $x$, $y$
  son enteros.

\item Podemos tomar decisiones con el comando if, else, y elseif
\begin{verbatim}
for x = 1:15
  if (mod(x, 2) == 0)
     fprintf('%d es par\n', x);
  elseif (mod(x, 3) == 0)
     fprintf('%d es impar y divisible por 3\n', x);
  else
     fprintf('%d es impar\n', x);
  end
end
\end{verbatim}

\end{enumerate}

\section{Ejemplo 1}

En esta secci\'{o}n vemos 
\begin{enumerate}
\item Como generar una lista de $51$ n\'{u}meros,
\item Imprimir en pantalla.
\item Graficar, colocando las leyendas correspondientes a título, 
      eje x, eje y.
\item Guardar el gráfico en un archivo.
\end{enumerate}
La idea es hacer una pr\'{a}ctica de la escritura de un programa (a\'{u}n cuando este no tenga otra utilidad).

Hay varias de formas de generar n\'{u}meros.
\begin{enumerate}
\item Con un ciclo for (Matlab)
\begin{verbatim}
for x = 1:51
 fprintf('%d\n', x);
end;
\end{verbatim}

\item Con un ciclo for (Octave)
\begin{verbatim}
for x = 1:51
 fprintf('%d\n', x);
endfor;
\end{verbatim}

\item Con un ciclo while (Matlab)
\begin{verbatim}
x = 1;
while (x < 51)
 fprintf('%d\n', x);
 x = x + 1;
endwhile
\end{verbatim}

\item Con un ciclo while (Octave)
\begin{verbatim}
x = 1;
while (x < 51)
 fprintf('%d\n', x);
 x = x + 1;
endwhile
\end{verbatim}

\item Generando un vector
\begin{verbatim}
x = 1:1:51
\end{verbatim}
\end{enumerate}

Para graficar los n\'{u}meros generados, pruebe con
\begin{verbatim}
x = 1:1:51 % los datos del eje x
y = rand(size(x)) % los datos del eje y, generados al azar
plot(x, y); 
title('n\'{u}meros generados por computadora');
xlabel('n\'{u}meros consecutivos');
ylabel('n\'{u}meros al azar');
\end{verbatim}

\textquestiondown Como puede variar el color y el tipo de las gr\'{a}ficas?

A continuaci\'{o}n se explica como se hace para guardar los resultados
en un archivo de texto.

\section{Ejemplo 2}

En esta secci\'{o}n veremos como escribir y leer archivos de texto.

\begin{enumerate}
\item Para leer y escribir archivos de texto, necesitamos primero abrir el
archivo. Para esto usamos una funci\'{o}n denominada ``fopen''.

\item ``fopen'' necesita como m\'{\i}nimo un nombre de archivo como
  argumento de entrada. Por eso escribimos \verb-fopen('nombre');-.

\item ``fopen'' devuelve un valor que tenemos que guardar en una
  variable. Por lo tanto, el comando se escribe
\begin{verbatim}
fid = fopen('nombre.txt');
\end{verbatim}
La variable fid guarda un n\'{u}mero que funciona como identificador
del archivo (``File IDentifier''). Si fopen falla, fid ser\'{a} igual
a $-1$.

\item Un archivo se puede abrir para leer, para escribir, o
  para agregar.

\item Cuando un archivo se abre para leer, el comando es
\begin{verbatim}
fid = fopen('archivo','r');
\end{verbatim}
Si el archivo no existe, el comando falla.

\item Cuando un archivo se abre para escribir, el comando es
\begin{verbatim}
fid = fopen('archivo','w');
\end{verbatim}
Si el archivo no existe, se crea.

\item Cuando un archivo se abre para agregar, el comando es
\begin{verbatim}
fid = fopen('archivo','a');
\end{verbatim}
Si el archivo no existe, se crea.

\item Existen otras opciones que puede ver escribiendo ``help
  fopen''. Por el momento, no las usaremos.

\item Para escribir texto, podemos usar ``fprintf'' de la siguiente
  forma
\begin{verbatim}
x = 2;
y = 3;
fprintf(fid, ``x = %d; y = %f'', x, y);
\end{verbatim}

\item Para leer, podemos usar ``fscanf'' de la siguiente
  forma
\begin{verbatim}
A = fscanf(fid, ``%d'');
\end{verbatim}

\item Una ventaja de Matlab es que muchas funciones est\'{a}n
  vectorizadas. Por lo tanto, escribiendo
\begin{verbatim}
[A,s] = fscanf(fid, ``%d'',[3,10]);
\end{verbatim}
podemos leer una matriz de 3 filas y 10 columnas (suponiendo que el
archivo contenga esa informaci\'{o}n.

\item Finalmente, tanto para lectura como para escritura, debemos
  cerrar los archivos con el comando ``fclose(fid)''.

\item Como ejemplo, veamos
\begin{verbatim}
% este programa escribe y lee 
% una matriz de 3 filas y 10 columnas

% abro archivo texto para escritura
esc = fopen('texto','w');

for x = 1:3
    for y = 1:10

        % escribo en archivo 
        % identificado con la
        % variable esc
        fprintf(esc, '%d ', x + y);
    end

    % escribo en archivo 
    % identificado con la
    % variable esc
    fprintf(esc,'\n');
end

% cierro archivo identificado con 
% la variable esc
fclose(esc);

% abro archivo texto para lectura
lec = fopen('texto','r');

% leo archivo identificado con 
% la variable lec
[A,s] = fscanf(lec,'%d',[3,10]);

% cierro archivo identificado con 
% la variable lec
fclose(lec);
\end{verbatim}

\item El archivo ``texto'' generado resulta ser
\begin{verbatim}
2 3 4 5 6 7 8 9 10 11 
3 4 5 6 7 8 9 10 11 12 
4 5 6 7 8 9 10 11 12 13 
\end{verbatim}

\item La variable A coincide con lo guardado en el archivo.

\item La variable s es igual a $30$, porque se leyeron correctamente
  $30$ n\'{u}meros enteros.

\end{enumerate}

\section{Ejemplo 3}

En esta secci\'{o}n vemos como 
\begin{enumerate}
\item Leer una tabla de datos con tres columnas a las que llamaremos a, b, y c.
\item Calcular las coordenadas $u = (a + b)/c, v = c / (a - b)$. Evitar problemas de división por cero.
\item Imprimir u y v.
\item Graficar los resultados u y v, colocando las leyendas correspondientes a t\'{\i}tulo, eje x, eje y.
\end{enumerate}

A continuaci\'{o}n vemos como realizar operaciones b\'{a}sicas con vectores y matrices.
\begin{enumerate}

\item Para generar un vector de $1001$ elementos, puedo escribir
\begin{verbatim}
x = 0:1:1000;
\end{verbatim}

\item Para el c\'{a}lculo del promedio debo sumar todos los elementos de
$x$. Explique porqu\'{e} funciona el siguiente truco:
\begin{verbatim}
u = ones(size(x));
suma = u' * x
\end{verbatim}
Para entenderlo, pruebe un vector $x$ con muy pocos elementos ($4$ o
$5$).

\item \textquestiondown C\'{o}mo calculan el promedio?

\item Usamos la constante eps de Matlab para protegernos contra la
  divisi\'{o}n  por cero.
\item La constante eps es el n\'{u}mero m\'{a}s peque\~{n}o distinto de $0$ que maneja
   Matlab.
\item.Si un divisor es $0$, le sumo la constante eps.
\item Para seleccionar toda la columna n de una matriz A, puedo escribir A(:,n).
\item Al escribir : en A(:, n), Matlab se encarga de usar un ciclo for para recu-
    perar todas las filas de la columna n de la matriz A.
\item Para seleccionar toda la fila n de una matriz A, puedo escribir A(n,:).
\item Al escribir : en A(n, :), Matlab se encarga de usar un ciclo for para recu-
    perar todas las columnas de la fila n de la matriz A.
\item Puedo verificar una condici\'{o}n sobre todo un vector sin escribir un ciclo.
    Supongamos que $x = [1, 2, -3, -4, -6, 7, 8, 9]$. Entonces $y = x > 0$ me
    permite obtener $y = [1, 1, 0, 0, 0, 1, 1, 1]$ donde los $1$ en las posiciones $1, 2,
    6, 7 y 8$ de y me indican que $x(1)$ es mayor que $0$, $x(2)$ es mayor que $0$,
    $x(6)$ es mayor que $0$, $x(7)$ es mayor que $0$ y $x(8)$ es mayor que $0$, y los $0$
    en las posiciones $3$, $4$ y $5$ de y me indican que $x(3) <= 0$, $x(4) <= 0$ y
    $x(5) <= 0$.
\item Supongamos que queremos dividir todos los elementos del vector $x =
    [1, 2, 3, 4]$, por los del vector $y = [2, 3, 4, 6]$ (ambos del mismo tama\~{n}o).
    Podemos dividir elemento a elemento a los vectores x e y escribiendo $x./y$.
\item Ejemplo
\begin{verbatim}
% este programa escribe y lee 
% una matriz de 10 filas y 3 columnas

% abro archivo texto para escritura
esc = fopen('texto','w');

for x = 1:10
    for y = 1:3

        % escribo en archivo 
        % identificado con la
        % variable esc
        fprintf(esc, '%d ', x + y);
    end

    % escribo en archivo 
    % identificado con la
    % variable esc
    fprintf(esc,'\n');
end

% cierro archivo identificado con 
% la variable esc
fclose(esc);

% abro archivo texto para lectura
lec = fopen('texto','r');

% leo archivo identificado con 
% la variable lec
[A,s] = fscanf(lec,'%d',[10,3]);

% cierro archivo identificado con 
% la variable lec
fclose(lec);

% obtengo las filas por separado
a = A(:,1);
b = A(:,2);
c = A(:,3);
% obtengo los denominadores 
% protegiendome de division por 0
ud = c + (abs(c) < eps) * eps;
vd = (a - b);
vd = vd + (abs(vd) < eps) * eps;
% observe el uso de ./ en lugar de /
u = (a + b) ./ ud; % division elemento a elemento
v = c ./ vd; % division elemento a elemento
\end{verbatim}

\end{enumerate}

